\documentclass[11pt]{article}
\usepackage[a4paper,margin=28mm]{geometry}
\usepackage[T1]{fontenc}
\usepackage{lmodern,amsmath,amssymb,amsthm,mathtools,microtype}
\usepackage{xurl,hyperref,enumitem}
\hypersetup{colorlinks=true,linkcolor=blue,citecolor=blue,urlcolor=blue,
 pdftitle={Vertex excess and the sharp exponent in simplex stability},
 pdfauthor={AI-assisted research manuscript for Entry 005}}
\setlength{\emergencystretch}{2em}
\numberwithin{equation}{section}
\newtheorem{theorem}{Theorem}[section]
\newtheorem{lemma}[theorem]{Lemma}
\newtheorem{proposition}[theorem]{Proposition}
\newtheorem{corollary}[theorem]{Corollary}
\theoremstyle{remark}\newtheorem{remark}[theorem]{Remark}
\newcommand{\R}{\mathbb R}
\newcommand{\E}{\mathbb E}
\newcommand{\conv}{\operatorname{conv}}
\newcommand{\vol}{\operatorname{Vol}}
\newcommand{\pyr}{\operatorname{pyr}}
\newcommand{\supp}{\operatorname{supp}}
\newcommand{\diag}{\operatorname{diag}}
\newcommand{\ip}[2]{\langle #1,#2\rangle}
\newcommand{\norm}[1]{\lVert #1\rVert}
\newcommand{\abs}[1]{\lvert #1\rvert}
\newcommand{\pathref}[1]{\nolinkurl{#1}}
\title{Vertex excess and the sharp exponent in simplex stability}
\author{Entry 005 research manuscript\\\small AI-assisted manuscript candidate}
\date{8 October 2026}
\begin{document}
\maketitle
\begin{abstract}
For a full-dimensional convex $d$-polytope with at most $d+1+r$ vertices,
we determine the sharp power of the projection-body pyramid defect that
controls containment in a dilation of every prescribed maximum-volume
inscribed simplex, centered at that same simplex's centroid. The exponent
is $1/\min(r,d-1)$. We give an explicit all-scale upper bound and a
truncation-and-pyramid family with an exact limiting lower constant.
The geometric step is a cone-volume estimate controlled by the number of
inner apex points: collapsing two base rays under projection improves the
usual cap exponent. A finite law on the actual facets supplies the facet
incidence needed to convert a vertex bound into this apex count. We include
the finite determinant assignment, enclosure, and prescribed-simplex
retention arguments so that the dependence on the original geometric defect
can be checked directly.
\end{abstract}

\noindent\textbf{Verification status.} The underlying traditional proof and
sharpness construction received a separate analytic review by an independent
model, with a PASS disposition. This is not professional peer review. The
present manuscript is an expanded exposition of that reviewed argument;
the new vertex-excess theorem has not been formalized in Lean. Earlier Lean
results are identified in Appendix~\ref{app:sources}, with no claim that an
older compiled main theorem certifies the new result. No priority or
first-discovery claim is made.

\setcounter{tocdepth}{1}
\tableofcontents
\section{The invariant and the main theorem}
\label{sec:main}
For a full-dimensional convex body $K\subset\R^d$, let $|K|$ denote its
Euclidean volume. For a \emph{unit} vector $u$, let $\pi_K(u)$ be the
$(d-1)$-dimensional volume of its orthogonal projection onto $u^\perp$.
We reserve the homogeneous notation
\[
 b_K(z)=\begin{cases}\norm z\pi_K(z/\norm z),&z\ne0,\\0,&z=0\end{cases}
\]
for nonunit vectors. The projection body $\Pi K$ has support function
$h_{\Pi K}=b_K$. If $L$ has dimension $p$, set
\[
 R(L)=\frac{|\Pi L|}{|L|^{p-1}},\qquad
 \pyr L=\conv\bigl((L,0)\cup\{(0,1)\}\bigr).
\]
We use exactly the projection-body pyramid invariant of the retained
Entry 005 project:
\begin{equation}\label{eq:actualdefect}
 a(K)=\left(\frac d{d+1}\right)^d
       \frac{R(\pyr K)}{R(K)}-1,
 \qquad e(K)=a(K)-\frac1{d+1}.
\end{equation}
For a full-dimensional simplex $S\subset K$ with centroid $g_S$, define
\begin{equation}\label{eq:excess}
 E(K,S)=\inf\{\tau\ge0:K\subset g_S+(1+\tau)(S-g_S)\}.
\end{equation}
The dilation factor itself is $1+E$. The simplex in this definition is
prescribed; neither the simplex nor its center is optimized in $E(K,S)$.

\begin{theorem}[Sharp vertex-excess exponent]\label{thm:main}
Let $d\ge3$ and $r\ge1$ be integers, and put
\[
 \begin{gathered}
 n=d+1,\qquad k=\min(r,d-1),\qquad H_d=3d^2(d+1)^3(d+2),\\
 T_{d,r}=\min\{64(d+1)r,16(d+1)^2\}.
 \end{gathered}
\]
For every full-dimensional convex $d$-polytope $K$ with at most $d+1+r$
vertices, and every prescribed maximum-volume inscribed simplex $S$,
\begin{equation}\label{eq:main}
 E(K,S)\le T_{d,r}H_d^{1/k}e(K)^{1/k}.
\end{equation}
The invariant $e(K)$ is nonnegative. The exponent $1/k$ is optimal:
there is an explicit family $(K_t,S_t)$ in this same class, with $S_t$
maximum and $e(K_t)\to0$ as $t\downarrow0$, such that
\begin{equation}\label{eq:sharp_limit}
 \lim_{t\downarrow0}\frac{E(K_t,S_t)}{e(K_t)^{1/k}}
 =n\left(\frac{n^2}{k(k+1)}\right)^{1/k}.
\end{equation}
Consequently no exponent $\beta>1/k$ admits a uniform bound near zero
defect, even with an arbitrary finite coefficient depending on $d,r$.
\end{theorem}

One extra vertex gives a linear defect estimate. Two extra vertices give
exponent $1/2$ in every dimension $d\ge3$. Saturation at the unrestricted
exponent $1/(d-1)$ occurs when the vertex allowance reaches $2d$.
The coefficient in \eqref{eq:main} is explicit but is not asserted to be
optimal. The case $r=0$ is trivial: a full-dimensional polytope with at most
$d+1$ vertices is a simplex. The theorem concerns $d\ge3$; the two-dimensional
body used in the sharpness construction is computed explicitly below.

The proof separates three geometric operations. First, a finite apex count
controls the volume missing from a cone and survives a ray-collapsing
projection. Second, actual facet normals yield an enclosing simplex with a
relative projection error of at most $H_de(K)$. Third, a near-vertex
condition for that enclosing simplex controls the originally prescribed
maximum, including maxima whose vertices are not vertices of $K$.

\subsection{Context and scope}
Mixed-volume polynomials, their monotonicity, the first mixed-volume formula,
and Minkowski's first inequality are classical inputs
\cite{SchneiderBook,BS}. They are not new ingredients of this work.
Simplex stability for other functionals is studied, for example, by
Schneider~\cite{SchneiderStability}; cone-volume stability for subspace
concentration and the $U$-functional is treated in~\cite{BH}. Projection-body
volume questions also have substantial earlier literature
\cite{Brannen,Saroglou}. These neighboring statements involve different
functionals or classes and are not used to infer a novelty claim here.
The source comparison accompanying this manuscript is focused rather than
exhaustive; in particular it does not certify the complete overlap with
all earlier projection-body results.

The cone count, projection collapse, and actual-facet specialization are
the added steps relative to the retained unrestricted proof. For transparency,
we also reproduce finite versions of its determinant assignment and
retention arguments. Thus the mathematical chain below can be read without
running the earlier formal source tree.

\section{Cone volume and a ray-collapsing projection}
\label{sec:cones}
\begin{lemma}[Finite apex count]\label{lem:cone}
Let $N,m$ be positive integers. Let $B$ be a compact convex
$(N-1)$-body in the affine hyperplane $\ell=1$ of $\R^N$, where $\ell$
is linear, and suppose $C=\conv(0,B)$ is full dimensional.
For $z_1,\ldots,z_m\in B$ and $0\le\rho\le s_j\le1$ with $\rho<1$, set
\[
 L=\conv(B,s_1z_1,\ldots,s_mz_m).
\]
Then
\begin{equation}\label{eq:cone}
 \frac{|L|}{|C|}\le1-\rho^{\min(m,N)}.
\end{equation}
The exponent and constant are sharp.
\end{lemma}
\begin{proof}
If $\rho=0$, the assertion is $L\subset C$. For $N=1$, the base is a point
and the interval length ratio is at most $1-\rho$. We may thus suppose
$N\ge2$ and $0<\rho<1$. Since
\[
 s_jz_j=\frac{1-s_j}{1-\rho}\rho z_j+
          \frac{s_j-\rho}{1-\rho}z_j,
\]
we have $L\subset\conv(B,\rho Z)$, where
$Z=\conv(z_1,\ldots,z_m)$ and $h=\dim Z\le\min(m-1,N-1)$.
Put $q=N-1$. At height $\ell=s\in[\rho,1]$, the latter hull has section
\begin{equation}\label{eq:section}
 aB+bZ,\qquad a=\frac{s-\rho}{1-\rho},\qquad
 b=\frac{\rho(1-s)}{1-\rho}.
\end{equation}
To interpret its $q$-volume, choose $b_0\in B$ and subtract $s b_0$.
Since $a+b=s$, the translated section is
$a(B-b_0)+b(Z-b_0)$ in $\ker\ell$.

Normalize mixed volume by $V(B[q])=|B|_q$. The classical volume polynomial is
\begin{equation}\label{eq:mv}
 |aB+bZ|_q=\sum_{j=0}^q\binom qj a^{q-j}b^j V(B[q-j],Z[j]).
\end{equation}
Translation invariance and monotonicity give
$0\le V(B[q-j],Z[j])\le |B|_q$. For $j>h$ these terms vanish:
translate $Z$ into an $h$-dimensional linear subspace; a product box bounds
$|B+tZ|_q$ by $O(t^h)$ for $t\ge1$, and nonnegativity of the polynomial
coefficients forces every higher coefficient to vanish. This also follows
from the essential-collection criterion in~\cite[Theorem 2.2]{BS}.

The substitution $s=\rho+(1-\rho)v$ gives
\[
 \binom qj\int_\rho^1a^{q-j}b^j\,ds
 =(1-\rho)\rho^j\binom qj\int_0^1v^{q-j}(1-v)^j\,dv
 =\frac{(1-\rho)\rho^j}{q+1}.
\]
The normal-height Jacobian is $1/\norm\ell$, while
$|C|=|B|_q/(N\norm\ell)$. Consequently
\[
 \frac{|L|}{|C|}\le(1-\rho)\sum_{j=0}^h\rho^j
 =1-\rho^{h+1}\le1-\rho^{\min(m,N)}.
\]
For sharpness when $m\le N$, take $C=\conv(0,e_1,\ldots,e_N)$,
$B=\conv(e_1,\ldots,e_N)$, and $z_j=e_j$. The omitted region has the
volume of the simplex with intercepts $\rho$ on the first $m$ coordinate
rays and one on the others, namely $\rho^m|C|$. For $m>N$, include all
$N$ base vertices and add redundant points; the omitted volume is
$\rho^N|C|$.
\end{proof}

\begin{remark}[Factorial convention]\label{rem:factorial}
With the diagonal normalization above, the coefficient of
$\lambda_1\cdots\lambda_q$ in
$|\lambda_1K_1+\cdots+\lambda_qK_q|_q$ is
$q!V(K_1,\ldots,K_q)$. Grouping repeated factors gives exactly the binomial
coefficient in \eqref{eq:mv}. The coefficient statement in
\cite[Theorem 2.1, version 2]{BS} appears to omit this factorial relative
to that paper's definition (1.1). We use its classical background and
multilinearity with the normalization stated here, not that coefficient
statement literally. At $q=2$ and $K_1=K_2=B$, the cross coefficient of
$(\lambda_1+\lambda_2)^2|B|_2$ is already $2|B|_2$.
\end{remark}

\begin{lemma}[Projection collapse]\label{lem:collapse}
Let $d\ge2$ and $r\ge1$ be integers. Let $P$ be a full-dimensional
$d$-simplex, let $K\subset P$ be a full-dimensional polytope, and fix a
vertex $v$ of $P$ with opposite facet $B$. Write $\lambda_v$ for its
$P$-barycentric coordinate and
$\rho=1-\max_{x\in K}\lambda_v(x)$. If at most $r+1$ vertices of $K$
lie outside $B$, some unit direction $u$ parallel to $B$ satisfies
\begin{equation}\label{eq:collapse}
 \frac{\pi_P(u)-\pi_K(u)}{\pi_P(u)}\ge\rho^{\min(r,d-1)}.
\end{equation}
\end{lemma}
\begin{proof}
For $\rho=0$ the result is immediate. Full dimensionality gives $\rho<1$.
Translate $v$ to zero and set $\ell=1-\lambda_v$, a linear functional
with $\ell(B)=1$. List the off-base vertices as $q_1,\ldots,q_j$,
where $1\le j\le r+1$. Uniquely,
\[
 q_i=s_i z_i,\qquad z_i\in B,\qquad \rho\le s_i<1.
\]
Thus $K\subset\conv(B,q_1,\ldots,q_j)$.
If $j\ge2$ and $z_1\ne z_2$, project orthogonally along $z_1-z_2$.
Their projected base points coincide. Of the projected $q_1,q_2$, retain
the one with the smaller $s_i$; the other lies between it and the common
projected base point. If $z_1=z_2$, the same redundancy exists before
projection, and any nonzero base-parallel direction can be chosen.
If $j=1$, choose any such direction and retain this single point.

In every case at most $m\le r$ inner apex points remain. The functional
$\ell$ descends through the projection. The projected $P$ is a genuine
$(d-1)$-cone with a $(d-2)$-dimensional base; all retained depths are at
least $\rho$. Lemma~\ref{lem:cone} gives missing fraction at least
$\rho^{\min(m,d-1)}\ge\rho^{\min(r,d-1)}$. Normalize the projection
direction to obtain the required unit $u$.
\end{proof}

\begin{corollary}[Actual facet incidence]\label{cor:caps}
Let $d\ge2$, $r\ge1$ be integers and $k=\min(r,d-1)$. Suppose $K\subset P$,
where $P$ is a $d$-simplex and $K$ has at most $d+1+r$ vertices. Assume
every facet hyperplane of $P$ is the hyperplane of an actual facet of $K$.
If, for $0\le\delta\le1$,
\[
 \pi_P(u)-\pi_K(u)\le\delta\pi_P(u)\qquad(\norm u=1),
\]
then every $P$-coordinate satisfies
\begin{equation}\label{eq:caps}
 \max_{x\in K}\lambda_i(x)\ge1-\delta^{1/k}.
\end{equation}
\end{corollary}
\begin{proof}
An actual facet of a full-dimensional $d$-polytope contains at least $d$
vertices. Thus at most $r+1$ vertices of $K$ lie outside any specified
facet of $P$. Apply Lemma~\ref{lem:collapse} at each opposite vertex.
\end{proof}
Mere support contact does not give this count: a supporting plane can touch
only one vertex. The next two sections construct $P$ from actual facets,
rather than assume that a limiting support law supplies this property.

\section{A finite determinant assignment theorem}
\label{sec:assignment}
We record a finite version of the retained volume-weighted assignment
argument, with its proof. For a real number $t$, write $t_+=\max(t,0)$.
All probability measures in this section have finite support, so no
integrability or limiting-law issue is present.

\begin{lemma}[Seminorm assignment]\label{lem:assignment}
Let $\mu$ be a centered probability measure with finite support
$\mathcal A\subset\R^d$ whose affine span is $\R^d$. Let $N$ be a seminorm
with $N(x)\le R_0$ on $\mathcal A$, where $R_0\ge0$. For independent
samples with law $\mu$, put
\[
 A=\E\abs{\det(X_1,\ldots,X_d)},\qquad
 B=\E\abs{\det((X_0,1),\ldots,(X_d,1))}.
\]
Then $B>0$, $D:=B-A\ge0$, and there are affinely independent anchors
$w_0,\ldots,w_d\in\mathcal A$ and an assignment $I:\mathcal A\to\{0,\ldots,d\}$
with
\begin{equation}\label{eq:assignment}
 h:=\E N(X-w_{I(X)})\le R_0(d+1)(d+2)\frac{D}{B}.
\end{equation}
One can choose $I(x)$ to be the least index of a largest barycentric
coordinate of $x$ in the anchor simplex.
\end{lemma}
\begin{proof}
For an ordered base $\mathbf x=(x_1,\ldots,x_d)$ define
\[
 F_{\mathbf x}(y)=\det((x_1,1),\ldots,(x_d,1),(y,1)),\quad
 P_{\mathbf x}=\E(F_{\mathbf x}(X))_+,\quad
 M_{\mathbf x}=\E(-F_{\mathbf x}(X))_+.
\]
Centering and affinity in the last column give
$\abs{\E F_{\mathbf x}(X)}=\abs{\det(x_1,\ldots,x_d)}$. Hence
\begin{equation}\label{eq:Dsign}
 D=2\E_{\mathbf X}\min(P_{\mathbf X},M_{\mathbf X})\ge0.
\end{equation}
For two test points set
\[
 \psi_{\mathbf x}(y,z)=
 \min((F_{\mathbf x}(y))_+,(-F_{\mathbf x}(z))_+)
 +\min((-F_{\mathbf x}(y))_+,(F_{\mathbf x}(z))_+).
\]
Each conditional expected minimum is at most
$\min(P_{\mathbf x},M_{\mathbf x})$. Therefore
\begin{equation}\label{eq:psibudget}
 \E_{\mathbf X,Y,Z}\psi_{\mathbf X}(Y,Z)\le D.
\end{equation}
These statements include singular bases.

For an anchor tuple $W=(w_0,\ldots,w_d)$, put
$v(W)=\abs{\det((w_0,1),\ldots,(w_d,1))}$. Form $G(W,x)$ by adding the
following witnesses, with any fixed order on each base:
\begin{enumerate}[label=(\roman*),leftmargin=2em]
\item for each $i$, the base is the other $d$ anchors and the test points
are $w_i,x$;
\item for each $i<j$, the base is $x$ and the other $d-1$ anchors, and the
test points are $w_i,w_j$.
\end{enumerate}
Reversing base orientation does not change $\psi$. There are
$M_d=(d+1)(d+2)/2$ summands. Each uses $d+2$ independent samples when
$W$ and $x$ are sampled, so, with $J(W)=\E_XG(W,X)$,
\[
 \E_WJ(W)\le M_dD,\qquad \E_Wv(W)=B>0.
\]
Positivity of $B$ follows from an affinely independent tuple in the support.
Weighting anchor tuples by $v(W)/B$ gives an anchor tuple with $v(W)>0$ and
\begin{equation}\label{eq:weightedchoice}
 \frac{J(W)}{v(W)}\le M_d\frac{D}{B}.
\end{equation}
Indeed its weighted mean is
$\E[J(W)\mathbf1_{\{v(W)>0\}}]/B\le M_dD/B$.
No assertion that $J$ vanishes at singular tuples is needed.

Fix this tuple, and write $\alpha_i(x)$ for its barycentric coordinates.
Determinant affinity gives, for witnesses of type (i),
$v(W)\min(1,(-\alpha_i)_+)$; those of type (ii) equal
\[
 v(W)\bigl[\min((\alpha_i)_+,(\alpha_j)_+)
           +\min((-\alpha_i)_+,(-\alpha_j)_+)\bigr].
\]
It follows that $G(W,x)\ge v(W)\phi(x)$, where
\[
 \phi(x)=\sum_i\min(1,(-\alpha_i(x))_+)
       +\sum_{i<j}\min((\alpha_i(x))_+,(\alpha_j(x))_+).
\]
Let $j=I(x)$ be a largest coordinate. Since $\sum_i\alpha_i=1$, we have
$\alpha_j>0$. Put $N_x=\sum_i(-\alpha_i)_+$ and
$U_x=\sum_{i\ne j}(\alpha_i)_+$. The support bounds yield both
\[
 N(x-w_j)\le2R_0(N_x+U_x),\qquad N(x-w_j)\le2R_0.
\]
The first follows by expanding $x-w_j=\sum_{i\ne j}\alpha_i(w_i-w_j)$;
the second uses $x,w_j\in\mathcal A$. The pairs containing the largest
coordinate give $U_x\le\sum_{i<l}\min((\alpha_i)_+,(\alpha_l)_+)$.
Also $\min(1,N_x)\le\sum_i\min(1,(-\alpha_i)_+)$. Thus
\[
 N(x-w_{I(x)})\le2R_0\min(1,N_x+U_x)
 \le2R_0\phi(x)\le2R_0\frac{G(W,x)}{v(W)}.
\]
Integrating and using \eqref{eq:weightedchoice} proves \eqref{eq:assignment}.
The reasoning includes $D=0$ without division by $D$.
\end{proof}

\section{From the actual facet law to an enclosing simplex}
\label{sec:enclosure}
\subsection{Raw facet identities and the original defect}
For a full-dimensional $p$-polytope $L$, let $u_i=s_i n_i$ be its
facet area-normal vectors and $b_i=s_i h_L(n_i)$, where $n_i$ is the
outward unit normal and $s_i>0$ is the facet area. Set
\begin{equation}\label{eq:rawsums}
 \mathcal H(L)=\sum_{|J|=p}\abs{\det(u_j:j\in J)},\qquad
 \mathcal L(L)=\sum_{|J|=p+1}\abs{\det((u_j,b_j):j\in J)}.
\end{equation}
The order chosen within a subset is immaterial. Translating $L$ replaces
$b_i$ by $b_i-\ip{u_i}{z}$ and leaves these sums unchanged.

Cauchy's formula gives $b_L(z)=\frac12\sum_i\abs{\ip{z}{u_i}}$.
Hence $\Pi L=\sum_i[-u_i/2,u_i/2]$, and the usual volume formula for a
zonotope gives $|\Pi L|=\mathcal H(L)$. Direct facet parametrization of the
unit-height pyramid gives its side pairs and bottom pair as
\begin{equation}\label{eq:pyrfacets}
 \left(\frac{(u_i,b_i)}p,\frac{b_i}p\right),\qquad ((0,-|L|),0),
\end{equation}
respectively. One way to check the factor $1/p$ is to integrate the
$(p-1)$-dimensional slices of each side face: its oriented area vector is
$(u_i,b_i)\int_0^1(1-t)^{p-1}dt=(u_i,b_i)/p$.
Expanding the corresponding minors gives
\begin{align}
 \mathcal H(\pyr L)&=p^{-(p+1)}\mathcal L(L)
                     +|L|p^{-p}\mathcal H(L),\label{eq:pyrH}\\
 \mathcal L(\pyr L)&=|L|p^{-(p+1)}\mathcal L(L),
 &|\pyr L|&=\frac{|L|}{p+1}.\label{eq:pyrL}
\end{align}
For the second identity, an all-side lifted minor is zero because its
last two coordinate rows coincide; minors using the bottom give the
stated expression. Substituting \eqref{eq:pyrH} into the definition of $a$
therefore proves the exact identity
\begin{equation}\label{eq:a_raw}
 a(L)=\frac{\mathcal L(L)}{p|L|\mathcal H(L)}.
\end{equation}
These computations concern the original projection bodies, not a substitute
invariant. The standard projection covariance
$\Pi(AL)=|\det A|A^{-T}\Pi L$ shows that $R$, and hence $a$ and $e$,
are affine invariant; affine maps of a base extend to its pyramid.
The quantity $E$ is affine invariant as well.

\subsection{Normalization and the finite cone law}
Affinely normalize the prescribed maximum $S$ so that its own centroid is
zero and its image contains the closed unit ball. Continue to write $K,S$
for the images. With $n=d+1$, maximum-volume replacement gives
\begin{equation}\label{eq:replacement}
 |\alpha_i^S(x)|\le1\quad(x\in K).
\end{equation}
Indeed replacing vertex $i$ by $x$ multiplies the absolute volume by
$|\alpha_i^S(x)|$. Since $\sum_i\alpha_i^S(x)=1$ and the centroid is zero,
$-x/d$ has $S$-coordinates $(1-\alpha_i^S(x))/d\ge0$. Thus
\begin{equation}\label{eq:asymmetry}
 -K\subset dS\subset dK,\qquad E(K,S)\le n.
\end{equation}
The second conclusion uses the exact centered-dilation description
\begin{equation}\label{eq:dilationcoords}
 g_S+(1+\tau)(S-g_S)=\{x:\alpha_i^S(x)\ge-\tau/n\text{ for all }i\}.
\end{equation}

For the actual irredundant facets $F$ of $K$, use the notation
$n_F,s_F,h_F=h_K(n_F)$ and put
\begin{equation}\label{eq:law}
 X_F=\frac{n_F}{h_F},\qquad p_F=\frac{s_Fh_F}{d|K|},\qquad
 \mu=\sum_F p_F\delta_{X_F}.
\end{equation}
Here $h_F\ge1$ because the unit ball is in $K$. The cone-volume identity
$\sum_Fs_Fh_F=d|K|$ and the surface-area balance $\sum_Fs_Fn_F=0$ show
that $\mu$ is a centered probability law. Its atoms are precisely actual
facet normals divided by their positive support numbers, with
$\norm{X_F}\le1$ and $h_K(X_F)=1$.
Boundedness of $K$ forces these atoms to span $\R^d$; centeredness with
positive weights then gives full affine span. In particular the horizontal
and lifted moments $A,B$ of Section~\ref{sec:assignment} are positive.

Cauchy's formula and centeredness give the exact brightness identities
\begin{equation}\label{eq:brightness}
 m_K(z):=\E(\ip{z}{X})_+=\frac{b_K(z)}{d|K|},\qquad
 f_K(z):=\E|\ip zX|=2m_K(z).
\end{equation}
Expanding the iid determinants by distinct facet subsets gives
\[
 A=\frac{d!\mathcal H(K)}{(d|K|)^d},\qquad
 B=\frac{(d+1)!\mathcal L(K)}{(d|K|)^{d+1}}.
\]
Together with \eqref{eq:a_raw}, this proves
\begin{equation}\label{eq:momentdefect}
 a(K)=\frac{B}{nA},\qquad
 \frac{B-A}{B}=\frac{ne(K)}{1+ne(K)}.
\end{equation}
Jensen's inequality, or \eqref{eq:Dsign}, proves $B\ge A$ and hence $e(K)\ge0$.

Let $N=h_{K-K}$, a continuous seminorm. At every atom,
\[
 N(X_F)=h_K(X_F)+h_K(-X_F)\le1+d=n
\]
by \eqref{eq:asymmetry}. Lemma~\ref{lem:assignment}, with $R_0=n$, gives
$d+1$ affinely independent \emph{actual atoms} $w_i$ and an assignment with
\begin{equation}\label{eq:hbudget}
 h=\E N(X-w_{I(X)})\le
 \frac{Ce(K)}{1+ne(K)},\qquad C=n^3(n+1).
\end{equation}
Thus every selected $w_i$ is an actual facet normal, a fact that is stronger
than merely knowing that an anchor belongs to an unspecified support law.

\subsection{Width and same-anchor enclosure}
For every $z,y\in\R^d$ one has
\begin{equation}\label{eq:width}
 |\ip zy|\le d\,m_K(z)N(y).
\end{equation}
To see this, first take unit $v$ and $c=\ip yv\ne0$. In coordinates
$x=q+tv$ with $q\in v^\perp$, the determinant-one shear
$(q,t)\mapsto(q,t+\ip yq/c)$ places $K$ in a cylinder with base the
projection of $K$ and height $N(y)/|c|$. Thus
$|c||K|\le N(y)\pi_K(v)$. Scaling $v$ to $z$ and using
\eqref{eq:brightness} proves \eqref{eq:width}; zero cases are immediate.

Let $p_i=\mu\{I=i\}$. Positive-part Lipschitz continuity gives
\begin{equation}\label{eq:positivecomp}
 \sum_i p_i(\ip z{w_i})_+\ge(1-dh)m_K(z).
\end{equation}
Assume $h<1/(2d)$. On the unit sphere, $m_K$ has a positive lower bound,
since $K$ contains the unit ball. Hence for every unit $z$ some anchor has
positive pairing with $z$, uniformly bounded away from zero. Support
separation puts zero in the interior of $\conv(w_0,\ldots,w_d)$.
Its polar
\begin{equation}\label{eq:outer}
 P=\{x:\ip{w_i}x\le1\text{ for }0\le i\le d\}
\end{equation}
is therefore a genuine $d$-simplex containing $K$. Every defining
inequality is an actual $K$-facet inequality. The corresponding full facet
of $K$ lies in a facet of $P$, giving exactly the hypothesis of
Corollary~\ref{cor:caps}.

For $z\in P$, the left side of \eqref{eq:positivecomp} is at most one;
therefore $m_K(z)\le2$. Equation~\eqref{eq:width} and support separation
then give the intrinsic containment
\begin{equation}\label{eq:intrinsiccontain}
 P\subset2d(K-K).
\end{equation}

\subsection{Probability correction and relative projection error}
Let $\lambda_i>0$ be the barycentric coordinates of zero in the anchor
simplex. Let $q_i$ be the vertex of $P$ opposite the facet with normal $w_i$.
The anchor-simplex coordinates satisfy
\begin{equation}\label{eq:polarcoords}
 \alpha_i^w(x)=\lambda_i(1-\ip{q_i}x).
\end{equation}
Both sides agree at each anchor. The actual cone probabilities of $P$
satisfy the same mass and mean equations as $\lambda_i$ and hence equal
$\lambda_i$. Put $c=\sum_i p_iw_i$. Since $\E X=0$,
\eqref{eq:intrinsiccontain} gives $|\ip{q_i}c|\le2dh$. Evaluating
\eqref{eq:polarcoords} at $c$ yields
\[
 p_i=\lambda_i(1-\ip{q_i}c)\ge(1-2dh)\lambda_i.
\]
Also \eqref{eq:width} and the absolute-value triangle inequality give
\[
 \sum_i p_i|\ip z{w_i}|\le f_K(z)+d m_K(z)h
 =(1+dh/2)f_K(z).
\]
Write $\theta=2dh<1$. The exact brightness formulas for $P$ and $K$ imply
\begin{equation}\label{eq:brightnesscomp}
 (1-\theta)\frac{\pi_P(u)}{|P|}
 \le(1+\theta/4)\frac{\pi_K(u)}{|K|}\qquad(\norm u=1).
\end{equation}

The finite first mixed-volume formula and Minkowski's first inequality give
\begin{equation}\label{eq:volumeroot}
 \left(\frac{|P|}{|K|}\right)^{1/d}
 \le\frac{V(K[d-1],P)}{|K|}
 =\E h_P(X).
\end{equation}
Indeed the last expression is
$\sum_Fs_Fh_P(n_F)/(d|K|)$. Since $h_P(w_i)=1$, subadditivity and
\eqref{eq:intrinsiccontain} bound it by $1+2dh=1+\theta$.
Combining \eqref{eq:brightnesscomp} and \eqref{eq:volumeroot},
\[
 \frac{\pi_K(u)}{\pi_P(u)}
 \ge\frac{1-\theta}{(1+\theta)^d(1+\theta/4)}.
\]
For factors in $[0,1]$, the deficit of their product is at most the sum of
their deficits. Thus the missing fraction above is at most
\begin{equation}\label{eq:scalar}
 \theta+d\theta+\theta/4
 =2d(d+5/4)h\le3d^2h\qquad(d\ge3).
\end{equation}
Here $1-(1+\theta)^{-d}\le d\theta$ follows from
$1-t^d\le d(1-t)$ for $0\le t\le1$.

\begin{proposition}[Finite actual-facet enclosure]\label{prop:enclosure}
For $d\ge3$, if $0\le e(K)\le H_d^{-1}$, the normalized $K$ is contained
in a simplex $P$ whose facet hyperplanes are actual $K$-facet hyperplanes,
and
\begin{equation}\label{eq:finalprojection}
 \pi_P(u)-\pi_K(u)\le H_de(K)\pi_P(u)\qquad(\norm u=1).
\end{equation}
\end{proposition}
\begin{proof}
Since $H_d=3d^2C$, \eqref{eq:hbudget} gives
$h\le1/(3d^2)<1/(2d)$ even at $e(K)=H_d^{-1}$. All strict enclosure gates
hold. Equations~\eqref{eq:scalar} and \eqref{eq:hbudget} give
$3d^2h\le H_de(K)$.
\end{proof}
Every step used the same finite actual facet law. No compact approximation,
subsequence, or limiting-law identification is needed for this proposition.

\section{Retention of every prescribed maximum simplex}
\label{sec:retention}
This section reproduces the retained retention estimates with the quantifier
on the original simplex explicit. The arguments use only containment and
near-vertex points, independently of how those points were obtained.

\begin{lemma}[Substochastic determinant matching]\label{lem:matching}
Let $U$ be a nonnegative $l\times l$ column-substochastic matrix, where
$l\ge1$. If $|\det U|\ge1-\eta$ and $0\le\eta<1/3$, then a permutation
of its columns makes every diagonal entry at least $1-\eta$.
\end{lemma}
\begin{proof}
Represent each column as the expectation of a random standard basis
vector, assigning any missing mass to the zero vector. Independent
selection represents the determinant as the expectation of a variable in
$\{-1,0,1\}$. If every column but column $j$ is sampled, each resulting
cofactor vector either vanishes or selects exactly one entry of column $j$
with a sign. Hence $|\det U|\le\max_iU_{ij}$. Choose an entry at least
$1-\eta$ in each column. If two chosen entries share a row, the independent
selections collide with probability at least $(1-\eta)^2$, so
\[
 |\det U|\le1-(1-\eta)^2\le2\eta<1-\eta,
\]
a contradiction. The selected rows are distinct.
\end{proof}

\begin{lemma}[Inverse bootstrap]\label{lem:bootstrap}
Let $K\subset P$, where $P$ is a $d$-simplex and $S$ is a prescribed
maximum-volume simplex in $K$. Suppose $q_i\in K$ have $P$-coordinates
$Q_{ii}\ge1-\rho$. Suppose that after relabeling the vertices of this same
$S$, their $P$-coordinate matrix $W$ satisfies $W_{ii}\ge1-\eta$, where
$0\le\eta\le1/8$ and $0\le\rho\le1/16$. Then
\begin{equation}\label{eq:bootstrap}
 E(K,S)\le8(d+1)\rho.
\end{equation}
\end{lemma}
\begin{proof}
The matrix $W$ is nonnegative column stochastic and invertible. Put
$B=W^{-1}$, $D=B-I$, $\beta=\max_{ij}|D_{ij}|$, and
$r_j=1-W_{jj}$. From $DW=I-W$ and column stochasticity,
\[
 D_{ij}=(\delta_{ij}-W_{ij})+
           \sum_{l\ne j}(D_{ij}-D_{il})W_{lj}.
\]
Consequently $\beta\le\eta+2\eta\beta$, whence
\begin{equation}\label{eq:firstinverse}
 \beta\le\frac\eta{1-2\eta}\le\frac16.
\end{equation}
The $ii$ entry of $BW=I$ gives
\[
 B_{ii}-1\ge\frac{(1-\beta)r_i}{1-r_i}\ge r_i/2.
\]
On the other hand, $q_i\in K$ and maximality of the same $S$ imply
$(BQ)_{ii}\le1$ by replacement. Since $Q_{ii}\ge1-\rho$,
$B_{ii}>0$, and $B_{ij}\ge-\beta$ for $i\ne j$,
\[
 B_{ii}-1\le\rho(B_{ii}+\beta)\le2\rho.
\]
Thus $r_i\le4\rho$ for every $i$. Reapply the inverse identity with this
improved column loss to obtain
\[
 \beta\le\frac{4\rho}{1-8\rho}\le8\rho.
\]
Every $S$-coordinate of each $P$-vertex is an entry of $B$ and is at least
$-8\rho$. Convex averaging gives the same lower bound throughout $P$.
Equation~\eqref{eq:dilationcoords} now yields \eqref{eq:bootstrap} about
$S$'s own centroid.
\end{proof}

\begin{proposition}[Two all-scale retention bounds]\label{prop:retention}
Let $d\ge1$ be an integer, $n=d+1$, and let $K\subset P$, where $K$ is a
full-dimensional compact convex body and $P$ a $d$-simplex. Suppose every
$P$-coordinate attains at least $1-\rho$ on $K$, where $\rho\ge0$.
For every prescribed maximum-volume inscribed simplex $S$,
\begin{equation}\label{eq:unrestrictedretention}
 E(K,S)\le16n^2\rho.
\end{equation}
If, additionally, $K$ is a polytope with at most $n+r$ vertices for an
integer $r\ge1$, then
\begin{equation}\label{eq:boundedretention}
 E(K,S)\le64nr\rho.
\end{equation}
The local estimate $E(K,S)\le8n\rho$ holds in the first case for
$\rho\le1/(16n)$ and in the second case for $\rho\le1/(64r)$.
\end{proposition}
\begin{proof}
The universal estimate $E(K,S)\le n$ follows from maximum-volume
replacement and \eqref{eq:dilationcoords}, without any normalization.
It handles the complementary large-$\rho$ branches of both bounds.

For the unrestricted local branch, choose $q_i\in K$ with $P$-coordinate
at least $1-\rho$ and let $Q$ be their column coordinate matrix. Independent
column selection gives
\[
 \det Q\ge2\prod_iQ_{ii}-1\ge1-2n\rho.
\]
Indeed the identity selection has probability $\prod_iQ_{ii}$ and every
other signed selection is at least $-1$. For $\rho\le1/(16n)$ the right
side is positive, so these points form a comparison simplex. Maximality of
$S$ gives $|\det W|\ge\det Q\ge1-\eta$ with $\eta=2n\rho\le1/8$.
Lemma~\ref{lem:matching} relabels only the vertices of this same $S$ so
that $W_{ii}\ge1-\eta$. Lemma~\ref{lem:bootstrap} proves the local bound.
Together with $E\le n$ for $\rho>1/(16n)$, this proves
\eqref{eq:unrestrictedretention}.

For the bounded-vertex branch assume $\rho\le1/(64r)$. Choose actual
vertices $q_i$ maximizing the $P$-coordinates. They are distinct because
$\rho<1/2$. Their nonnegative column-stochastic matrix satisfies
\begin{equation}\label{eq:Qinverse}
 \norm{Q-I}_1\le2\rho,
 \qquad \norm{Q^{-1}}_1\le(1-2\rho)^{-1}\le1+4\rho.
\end{equation}
Here $\norm{\cdot}_1$ is the induced column-sum matrix norm and
$\rho\le1/4$. Thus $T=\conv(q_i)$ is a simplex in $K$.

First let $S_v$ be any maximum simplex with vertices among those of $K$.
Such a maximum exists by multiaffinity of the determinant. It omits at most
$r$ of the $q_i$ because there are at most $r$ other vertices. Let $R$ be
the omitted index set, of size $l\le\min(r,n)$. If $l=0$, then $S_v=T$.
Otherwise label its retained $q_i$ by their own indices and put the remaining
vertices in columns $R$. If $W$ is its coordinate matrix and $C_0=Q^{-1}W$,
then $(C_0)_i=e_i$ for $i\notin R$, so
\[
 |\det (C_0)_{R,R}|=|\det C_0|=\frac{|S_v|}{|T|}\ge1.
\]
The restricted matrix $U=W_{R,R}$ is column substochastic. Each column of
$(C_0)_{R,R}$ has $\ell^1$ norm at most $1+4\rho$; by
$W-C_0=(Q-I)C_0$ its distance from the corresponding $U$-column is at most
$2\rho/(1-2\rho)\le4\rho$. Telescoping the determinant gives
\[
 \bigl|\det U-\det(C_0)_{R,R}\bigr|
 \le4l\rho(1+4\rho)^{l-1}\le8l\rho.
\]
For the last inequality, $4l\rho\le1/16$ and
$(1+4\rho)^{l-1}\le(1-4l\rho)^{-1}<2$, by the binomial expansion.
Hence $|\det U|\ge1-\eta$ with $\eta=8r\rho\le1/8$.
Lemma~\ref{lem:matching} matches the omitted indices, while the retained
$q_i$ already have diagonal coordinates at least $1-\rho\ge1-\eta$.
Every maximum vertex simplex therefore has exactly one vertex in each of
the disjoint convex caps
\begin{equation}\label{eq:disjointcaps}
 \mathcal C_i=\{x\in K:\lambda_i^P(x)\ge1-\eta\}.
\end{equation}
This conclusion also covers $l=0$.

Now return to the prescribed maximum $S$, whose vertices need not be
vertices of $K$. Express each vertex of $S$ as a convex combination of
$K$-vertices and sample the combinations independently. Its determinant is
the expectation of the sampled determinants. Every sampled absolute
volume is at most $|S|$. Equality in the expectation and triangle
inequality forces every positive-probability sampled tuple to be a maximum
vertex simplex, with one point in each cap \eqref{eq:disjointcaps}.
For each $i$, some sampled column must lie in $\mathcal C_i$ with probability
one; otherwise independence would give positive probability that all
columns avoid it. Distinct caps require distinct such columns. There are
$n$ columns and $n$ caps, so after relabeling, each actual vertex of this
same $S$ lies in its corresponding cap by convexity.

Thus $W_{ii}\ge1-\eta$ for the prescribed $S$, and
Lemma~\ref{lem:bootstrap} proves $E(K,S)\le8n\rho$ locally.
For $\rho>1/(64r)$ the universal estimate gives $E(K,S)<64nr\rho$.
No change of simplex carrier or centroid has occurred.
\end{proof}

\section{Assembly of the upper estimate}
\label{sec:assembly}
We prove \eqref{eq:main}, leaving the sharpness statement to the next
section. Fix any of the maximum simplices $S$ specified in
Theorem~\ref{thm:main}, and normalize that same simplex as in
Section~\ref{sec:enclosure}. The affine map preserves vertex count,
maximality, $e(K)$, $E(K,S)$, and the correspondence of centroids.

Suppose first that $0\le e(K)\le H_d^{-1}$. Proposition~\ref{prop:enclosure}
provides an actual-facet-supported simplex $P$ with relative projection
error $\delta=H_de(K)\in[0,1]$. Since $K$ has at most $n+r$ vertices,
Corollary~\ref{cor:caps} gives one point of $K$ in each cap of coordinate
loss at most
\[
 \rho=\delta^{1/k},\qquad k=\min(r,d-1).
\]
Both conclusions of Proposition~\ref{prop:retention} concern the same
$K,S,\rho$. Taking their minimum yields
\[
 E(K,S)\le\min(64nr,16n^2)\rho
 =T_{d,r}H_d^{1/k}e(K)^{1/k}.
\]
At $e(K)=0$, this gives $\rho=0$ and $E(K,S)=0$ without division by
any vanishing parameter. At $e(K)=H_d^{-1}$ the strict enclosure gate
remains strict, as checked in Proposition~\ref{prop:enclosure}.

For $e(K)>H_d^{-1}$, we have $(H_de(K))^{1/k}>1$ and
$T_{d,r}\ge n$. Therefore the universal replacement bound $E(K,S)\le n$
proves the same estimate. Undoing the affine map yields \eqref{eq:main}
for the originally prescribed maximum and its original centroid.

The exponent is supplied by the finite apex count; the constants $64nr$
and $16n^2$ come from retention. Their minimum does not select different
maxima or different centers. No classification of all polytopes with two
extra vertices, stronger all-minor stochastic conjecture, or forest
conjecture enters this argument.

\section{The exact truncation and pyramid obstruction}
\label{sec:sharpness}
Put $k=\min(r,d-1)$ and $m=k+1$, so $2\le m\le d$. For $0<t<1$ define
\begin{equation}\label{eq:core}
 C_t=\left\{x\in\R^m:x_i\ge0,\quad t\le\sum_i x_i\le1\right\}
     =\conv\{e_1,\ldots,e_m,te_1,\ldots,te_m\}.
\end{equation}
This is a simplex truncated parallel to the facet opposite zero.

\begin{proposition}[Exact core calculation]\label{prop:core}
The body $C_t$ has $2m$ vertices, and
$T_t=\conv(e_1,\ldots,e_m,te_1)$ is a maximum-volume inscribed simplex.
With the original centroid of $T_t$,
\begin{equation}\label{eq:coreE}
 |T_t|=\frac{1-t}{m!},\qquad E(C_t,T_t)=(m+1)t.
\end{equation}
Its actual defect is
\begin{equation}\label{eq:coree}
 e_m(C_t)=
 \frac{t^{m-1}[m(m-1)-(m+1)(m-2)t-2t^m]}
 {(m+1)(1-t^m)[m+1+(m-1)t^{m-1}]}.
\end{equation}
In particular,
\begin{equation}\label{eq:corelimit}
 \lim_{t\downarrow0}\frac{e_m(C_t)}{t^{m-1}}
 =\frac{m(m-1)}{(m+1)^2}>0.
\end{equation}
\end{proposition}
\begin{proof}
Every full-dimensional vertex simplex must use all $m$ coordinate rays;
otherwise it lies in a coordinate hyperplane. An $(m+1)$-tuple consequently
has exactly one doubled ray. Its determinant magnitude is
$(1-t)t^j$ for some $j\ge0$: the doubled ray contributes $1-t$ and each
other selected low vertex contributes a factor $t$. Hence the maximum
vertex-simplex volume is $(1-t)/m!$, attained by $T_t$. Expressing arbitrary
inscribed points as convex combinations of the vertices and using
multiaffinity and the triangle inequality proves the same bound for all
inscribed simplices.

Write $s=\sum_i x_i$. Relative to the ordered vertices
$(te_1,e_1,e_2,\ldots,e_m)$ of $T_t$, the coordinates are
\[
 \beta_0=\frac{1-s}{1-t},\qquad
 \beta_1=x_1-t\beta_0,\qquad \beta_i=x_i\quad(2\le i\le m).
\]
On $C_t$, $0\le\beta_0\le1$ and all coordinates are at least $-t$.
At $te_2$, the coordinate $\beta_1$ equals $-t$. The centered-dilation
formula \eqref{eq:dilationcoords} proves \eqref{eq:coreE}.

For the defect put $c=1/(m-1)!$, $q=t^{m-1}$, and $b=1-q$.
The actual facet area-normal and support pairs are
\[
 (-cb e_i,0)\quad(1\le i\le m),\qquad
 (c\mathbf1,c),\qquad(-cq\mathbf1,-ct^m).
\]
The volume is $(1-t^m)/m!$. The horizontal minors consist of the
all-coordinate subset, the $m$ subsets using the top facet, and the $m$
subsets using the bottom facet; subsets using both parallel top and bottom
normals vanish. Therefore
\begin{equation}\label{eq:coreH}
 \mathcal H(C_t)=c^m b^{m-1}[m+1+(m-1)q].
\end{equation}
For the lifted minors the all-coordinate-plus-top term is $c^{m+1}b^m$,
the all-coordinate-plus-bottom term is $c^{m+1}b^m t^m$, and each of the
$m$ terms with top, bottom, and $m-1$ coordinate facets is
$c^{m+1}b^{m-1}(q-t^m)$. Thus
\begin{equation}\label{eq:coreL}
 \mathcal L(C_t)=c^{m+1}b^{m-1}
       [1+(m-1)q-(m-1)t^m-qt^m].
\end{equation}
By \eqref{eq:a_raw},
\[
 a_m(C_t)=
 \frac{1+(m-1)q-(m-1)t^m-qt^m}
 {(1-t^m)[m+1+(m-1)q]}.
\]
Subtracting $1/(m+1)$ gives \eqref{eq:coree}. Its numerator bracket equals
$(m+1)(m-2)(1-t)+2(1-t^m)>0$ for every $m\ge2$.
The limit \eqref{eq:corelimit} follows.
\end{proof}

\begin{remark}[The two-dimensional core]
The computation explicitly includes $m=2$; it is not an invocation of a
statement restricted to dimensions at least three. In this case
\begin{equation}\label{eq:twoD}
 e_2(C_t)=\frac{2t}{3(3+t)},\qquad
 3-\frac1{a_2(C_t)}=\frac{2t}{1+t}.
\end{equation}
This is the core required for the entire $r=1$ layer of the main theorem.
\end{remark}

\begin{lemma}[Pyramid identities]\label{lem:pyramid}
For a full-dimensional $p$-polytope $L$ with $a_p(L)>0$,
\begin{equation}\label{eq:apyr}
 a_{p+1}(\pyr L)=\frac{a_p(L)}{1+a_p(L)}.
\end{equation}
Consequently
\begin{equation}\label{eq:Deltainvariant}
 \Delta(L):=p+1-\frac1{a_p(L)}
\end{equation}
is unchanged by taking a pyramid. If $T$ is a maximum simplex in $L$,
then $\pyr T$ is a maximum simplex in $\pyr L$. If all $T$-coordinates
on $L$ are at least $-t$, and one coordinate attains $-t$, then the
same is true for the $\pyr T$-coordinates on $\pyr L$.
\end{lemma}
\begin{proof}
Substitute \eqref{eq:pyrH}--\eqref{eq:pyrL} into \eqref{eq:a_raw} to get
\[
 a_{p+1}(\pyr L)
 =\frac{\mathcal L(L)}{\mathcal L(L)+p|L|\mathcal H(L)}
 =\frac{a_p(L)}{1+a_p(L)}.
\]
Thus $1/a_{p+1}=1+1/a_p$, proving \eqref{eq:Deltainvariant}.
Every full-dimensional vertex simplex in a pyramid includes its unique
apex, and its base is a full-dimensional vertex simplex of $L$. Comparing
base volumes proves the maximum statement for vertex simplices;
multiaffinity extends it to arbitrary inscribed simplices. A point in the
pyramid can be written as $((1-s)x,s)$ with $x\in L$, $0\le s\le1$.
Its old coordinates are $(1-s)\alpha_i^T(x)$ and its new apex coordinate
is $s$. They are bounded below by $-t$, and the old base still realizes
that lower value.
\end{proof}

\begin{proof}[Proof of sharpness in Theorem~\ref{thm:main}]
Set
\[
 K_t=\pyr^{\,d-m}C_t,\qquad S_t=\pyr^{\,d-m}T_t.
\]
These are actual $d$-polytopes and actual prescribed maxima. Their vertex
count is $2m+(d-m)=d+1+k\le d+1+r$. Iterating
Lemma~\ref{lem:pyramid} and using \eqref{eq:coreE} gives
\begin{equation}\label{eq:ambientE}
 E(K_t,S_t)=nt,\qquad n=d+1.
\end{equation}
From \eqref{eq:corelimit} and the identity
\[
 \Delta(C_t)=\frac{(m+1)^2e_m(C_t)}{1+(m+1)e_m(C_t)},
\]
we obtain $\Delta(C_t)/t^k\to k(k+1)$. In ambient dimension $d$, pyramid
invariance and $a_d=1/(n-\Delta)$ yield exactly
\begin{equation}\label{eq:ambiente}
 e_d(K_t)=\frac{\Delta(C_t)}{n[n-\Delta(C_t)]},\qquad
 \frac{e_d(K_t)}{t^k}\longrightarrow\frac{k(k+1)}{n^2}.
\end{equation}
Together with \eqref{eq:ambientE}, this is \eqref{eq:sharp_limit}.
For $\beta>1/k$, the ratio $E(K_t,S_t)/e_d(K_t)^\beta$ is asymptotic to a
positive constant times $t^{1-k\beta}$, which diverges. Hence no finite
coefficient depending on $d,r$ can support that stronger exponent near
zero defect.
\end{proof}

\section{Exact sharpness of the geometric projection step}
\label{sec:geometricsharp}
The preceding family proves sharpness for the original invariant.
Here we separately verify sharpness of the existence-of-direction bound
in Lemma~\ref{lem:collapse}. Equality in just one direction would not
suffice for this purpose; an upper bound in every direction is needed.

Fix integers $1\le r\le d-1$, put $m=r+1$, and let $0<\rho<1$. Set
\begin{align*}
 P&=\conv(0,e_1,\ldots,e_d),\\
 K&=\conv(e_1,\ldots,e_d,\rho e_1,\ldots,\rho e_m)\\
  &=\left\{x\ge0:\sum_i x_i\le1,\quad
       \rho^{-1}\sum_{i\le m}x_i+\sum_{i>m}x_i\ge1\right\}.
\end{align*}
The barycentric loss at zero is $\rho$. This body has $d+m=d+1+r$
vertices and is an iterated pyramid over the $m$-dimensional equal
truncation. Every facet plane of $P$ supports a full facet of $K$.

With $c=1/(d-1)!$, the actual facet area-normal vectors of $K$ are
\begin{align*}
 &-c(1-\rho^{m-1})e_i &&(i\le m),\\
 &-c(1-\rho^m)e_i &&(i>m),\\
 &c\mathbf1,\qquad
 -c(\underbrace{\rho^{m-1},\ldots,\rho^{m-1}}_{m},
       \underbrace{\rho^m,\ldots,\rho^m}_{d-m}).
\end{align*}
For a coordinate facet the removed simplex has the corresponding product
of intercept scales; the last vector is obtained from the intercepts of
the new cut facet. These vectors sum to zero, as required.

For nonzero $z\in\R^d$ write
\[
 A=\sum_{i\le m}z_i,\quad B=\sum_{i>m}z_i,\quad
 A_* =\sum_{i\le m}|z_i|,\quad B_* =\sum_{i>m}|z_i|.
\]
Cauchy's formula gives the exact \emph{homogeneous} brightness ratio
\begin{equation}\label{eq:sharpbrightness}
 \frac{b_P(z)-b_K(z)}{b_P(z)}
 =\rho^{m-1}\frac{A_*+\rho B_*-|A+\rho B|}
                         {A_*+B_*+|A+B|}.
\end{equation}
The denominator is positive. The numerator is nonnegative by the triangle
inequality and is at most the denominator. Therefore the relative deficit
is at most $\rho^{m-1}=\rho^r$ in every direction. For
$z=e_1-e_2$, the ratio equals $\rho^r$, and that direction is parallel to
the opposite base facet of $P$. Upon normalizing $z$, we conclude
\begin{equation}\label{eq:maxsharp}
 \max_{\norm u=1}\frac{\pi_P(u)-\pi_K(u)}{\pi_P(u)}=\rho^r.
\end{equation}
The same maximum is attained among base-parallel unit directions.
Thus both the exponent and the constant one in
Lemma~\ref{lem:collapse} are sharp throughout $1\le r\le d-1$.
The distinction between $b_K(z)$ and $\pi_K(u)$ avoids any convention for
projection notation at nonunit vectors.

\appendix
\section{Source correspondence and verification boundary}
\label{app:sources}
This manuscript expands the frozen proof
\pathref{research_math/vertex_excess_sharp_exponent_20261008/SHARP_VERTEX_EXCESS_EXPONENT.md}
and its geometric-sharpness addendum. The frozen main has SHA-256
\pathref{8da4c012d4099465b4702ce4091ba4d1e2cc7862c12e8eb55cbd7467414e0f05}.
The separate traditional review is
\pathref{research_math/vertex_excess_exponent_independent_audit_20261008/REFEREE_REPORT.md},
SHA-256
\pathref{d638e1ebf4505f480d99c08c21a2012fb36f0e021cd86d828d142effd5df283b}.
The supplied review checks the main exponent, the exact actual-defect
constant, the arbitrary-prescribed-maximum quantifier, both defect
endpoints, and the lower family. It is an independent model's analytic
review, not a human referee report or a new Lean replay.

\subsection{Retained arguments reproduced here}
The volume-weighted determinant selection is retained from
\pathref{notes/quadratic-dimensional-simplex-stability/imports/weighted_anchors.md}
in the project publication tree. Section~\ref{sec:assignment} supplies its
finite-support seminorm proof. The bounded-vertex retention argument is
retained from
\pathref{research_math/linear_retention_investigation_20261008/BOUNDED_VERTEX_EXCESS_RETENTION.md}.
Section~\ref{sec:retention} includes both that proof and the unrestricted
retention proof. The exact truncation calculation is retained from
\pathref{notes/simplex-truncation-stability/proof.tex} in the publication
tree; Section~\ref{sec:sharpness} includes the facet calculation, the
$m=2$ case, the maximum-simplex proof, and the pyramid identities.
The package's \pathref{SOURCE_MAP.md} records the precise local source
paths, hashes, and section-level correspondence, and includes frozen
copies of the main input documents.

\subsection{Lean references refer to retained statements only}
For the earlier formal source, the module filenames and theorem namespaces
must be distinguished. In particular, the retained interfaces include
\begin{itemize}[leftmargin=1.5em]
\item \pathref{Entry005.finite_halfspace_entryA_lifted_moment}, in
\pathref{Entry005/PyramidEntryDefect.lean};
\item \pathref{Entry005.finite_halfspace_defect_over_first_moment}, in
\pathref{Entry005/PyramidMomentDefect.lean};
\item \pathref{Entry005.unit_ball_seminorm_first_moment_assignment}, in
\pathref{Entry005/SeminormFirstMoment.lean}.
\end{itemize}
Actual finite-law probability, centering, ball support, and brightness
interfaces are in \pathref{Entry005/FiniteHalfspaceConeLaw.lean};
finite-support assignment interfaces are in
\pathref{Entry005/FiniteBodyWeightedAssignment.lean}; the earlier
unrestricted prescribed-simplex interface is in
\pathref{Entry005/RelativeSimplexRetention.lean}.
These are references to retained code, not new formalization claims for
this manuscript. In particular, \pathref{PyramidEntryDefect} and
\pathref{PyramidMomentDefect} are module names, not the theorem namespace.

The older compact-limit main theorem does not by itself provide the
actual-facet incidence used in Corollary~\ref{cor:caps}. Here that incidence
is proved directly by choosing anchors from the finite law
\eqref{eq:law}. The new cone/prismoid estimate, ray-collapse/counting step,
finite-law specialization with this incidence, and final exponent assembly
still need corresponding Lean formalization. An earlier successful
\pathref{compiledMain} or other main build cannot certify those additions.

\subsection{Editorial clarifications implemented}
The positive-integer restrictions in Lemma~\ref{lem:cone}, unit projection
convention and homogeneous brightness notation, explicit two-dimensional
core, actual Lean namespace references, and mixed-volume factorial
convention implement all nonblocking clarifications requested by the
supplied review. The finite actual-facet chain is expanded in full in
Section~\ref{sec:enclosure}. Standard mixed-volume inputs remain credited
as classical, and no strict-monotonicity theorem from~\cite{BS} is needed.

The theorem is a mathematical manuscript candidate. Neither the focused
literature check nor the separate model review establishes priority,
journal suitability, or acceptance by professional mathematicians.

\clearpage
\begin{thebibliography}{9}
\bibitem{BS}
F.~Bihan and I.~Soprunov,
\emph{Criteria for strict monotonicity of the mixed volume of convex polytopes},
arXiv:1702.07676v2 (2017).
\url{https://arxiv.org/abs/1702.07676v2}.
\bibitem{SchneiderBook}
R.~Schneider, \emph{Convex Bodies: The Brunn--Minkowski Theory},
second expanded edition, Encyclopedia of Mathematics and its Applications
151, Cambridge University Press, 2014.
\bibitem{SchneiderStability}
R.~Schneider, \emph{Stability for some extremal properties of the simplex},
Journal of Geometry \textbf{96} (2009), 135--148.
\url{https://doi.org/10.1007/s00022-010-0028-0}.
\bibitem{BH}
K.~J.~B\"or\"oczky and M.~Henk, \emph{Cone-volume measure and stability},
Advances in Mathematics \textbf{306} (2017), 24--50.
\url{https://arxiv.org/abs/1407.7272}.
\bibitem{Saroglou}
C.~Saroglou, \emph{Volumes of projection bodies of some classes of convex bodies},
Mathematika \textbf{57} (2011), 329--353.
\url{https://doi.org/10.1112/S0025579311001860}.
\bibitem{Brannen}
N.~S.~Brannen, \emph{Volumes of projection bodies},
Mathematika \textbf{43} (1996), 255--264.
\url{https://doi.org/10.1112/S002557930001175X}.
\end{thebibliography}
\end{document}
